\paragraph{Idea of the proof.}
The proof is in Appendix~\ref{app:bulk}. If we allow every jump size
$d\ge1$, the characteristic function of the compound Poisson law is
$e^{x\chi(e^{i\theta})}$ with $\chi(s)=(1-s)\log(1-s)/s$, and the jumps
$d\ge N$ change it little. For small $\theta$ we have
$\chi(e^{i\theta})=-i\theta\log(-i\theta)+O(\theta^2\lvert\log\lvert\theta\rvert\rvert)$.
With $\theta=t/x$ the exponent becomes $B(t)+it\log x$ plus an error of
order $t^2\log x/x$. The term $it\log x$ is the centering at $x\log x$. Both
characteristic functions decay like $e^{-c\lvert t\rvert}$, so Fourier
inversion turns this into a bound for each bin, and keeping the constants gives
$889$ and $165$. Part (b) then follows from Theorem~\ref{thm:cp}.



As an independent check, \texttt{certify\_center.py} encloses the crossing
and verifies $P_N(a)>P_N(h)$ at every noncentral interior bin for each even
$10\le N\le2000$. The $N=8$ check is in Lean. These center checks are no
longer needed for Corollary~\ref{cor:unique-min}, but the same crossing
brackets still supply the finite endpoint-dip input.



\begin{remark}[what is still open]\label{rem:convex-open}
The remaining gap concerns the prescribed convexity window, not the minimum
or the two-hump shape of Corollary~\ref{cor:unique-min}.
We do not compute the constant $N_1$ of Theorem~\ref{thm:convex-large}. The constants of
Theorem~\ref{thm:landau} are explicit, but the proof uses them only through a limit. Even if
we followed them, Step~3 needs $f_L(\frac32\lambda_{\mathrm L}+\frac12\log x)<\frac12f_L(\lambda_{\mathrm L})$ up
to the error $\Theta_N$. Numerically $f_L$ falls to half of its maximum at
$\lambda=2.432$. So even without the error $\Theta_N$, Step~3 needs
$\frac32\lambda_{\mathrm L}+\frac12\log x>2.432$, that is, $x>252.8$. Since
$x_N+\log x_N$ is close to $2\log N-\log8$ (Theorem~\ref{thm:erw-crossing}), this means that $N$
is above about $3.5\cdot10^{56}$. So the $N$ between $N_{\rm c}$ and $N_1$ (or the threshold $N_1'$ of Proposition~S3.1(c)
of Online Resource~1 for odd $N$) are open, and no
computation can reach these constants as the proof stands. By Corollary~\ref{cor:one-point}
(and Proposition~S3.1(b) of Online Resource~1 for odd $N$) the whole range would
be closed by an effective proof of the single inequality $D_N(a_1)>0$ at one point
$m^*<a_1\le2m^*+1$. Near $2m^*$ the Landau heuristic suggests second differences of size
$3/[x^3(\log x)^4]$, since $f_L''(\lambda)\sim6/\lambda^4$ and
$\lambda\approx\log x$ there. The pointwise error in Theorem~\ref{thm:landau}
is much larger, so it does not prove the needed sign.
\end{remark}

The density has a unique mode $\lambda_{\mathrm L}$ with
$-0.22278299<\lambda_{\mathrm L}<-0.22278297$
(Lemma~\ref{lem:landau-mode}, computer-assisted).


So at the crossing the law of $A_N$ has a hump near each end, about
$x\log x-0.2228\,x$ bins from it. The final position is $2A_N-N$, so in
physical distance each hump is twice as far from $\pm N$. For the shape of the law between
the humps let $D_N(a)=P_N(a+1)-2P_N(a)+P_N(a-1)$ be its second difference,
and for a mode $a^*\ge N/2$ of $P_N$ put $m^*=N-a^*$.
